\chapter{Sheaves}

\section{Presheaves and Sheaves}

\begin{definition}
Let \(X\) be a topological space and let \(\mathsf{C}\) be one of the
categories of sets, abelian groups, rings, or modules over a fixed ring.  A
\emph{presheaf} \(\FF\) on \(X\) with values in \(\mathsf{C}\) consists of an
object \(\FF(U)\) for every open set \(U\subset X\), and a restriction morphism
\[
    \rho^U_V:\FF(U)\longrightarrow \FF(V)
\]
for every inclusion \(V\subset U\), such that \(\rho^U_U=\id\) and
\(\rho^V_W\circ \rho^U_V=\rho^U_W\) whenever \(W\subset V\subset U\).  We write
\(s|_V\) for \(\rho^U_V(s)\).
\end{definition}

\begin{definition}
A presheaf \(\FF\) is a \emph{sheaf} if for every open set \(U\subset X\), every
open covering \(U=\bigcup_i U_i\), and every family \(s_i\in \FF(U_i)\)
satisfying
\[
    s_i|_{U_i\cap U_j}=s_j|_{U_i\cap U_j}
    \quad\text{for all }i,j,
\]
there exists a unique section \(s\in \FF(U)\) such that \(s|_{U_i}=s_i\) for
all \(i\).
\end{definition}

\begin{remark}
The sheaf condition has two parts.  Uniqueness says that sections are determined
locally.  Existence says that compatible local sections glue.
\end{remark}

\begin{example}
For a topological space \(Y\), the assignment
\[
    U\longmapsto \{ \text{continuous maps } U\to Y\}
\]
is a sheaf of sets on \(X\).  The sheaf property is precisely the local nature
of continuity and equality of functions.
\end{example}

\begin{example}
If \(A\) is an abelian group, the \emph{constant presheaf}
\(U\mapsto A\) with identity restrictions is usually not a sheaf when \(X\) is
disconnected.  Its sheafification is the sheaf of locally constant
\(A\)-valued functions.
\end{example}

\begin{definition}
A morphism of presheaves \(\varphi:\FF\to \GG\) is a family of morphisms
\(\varphi_U:\FF(U)\to \GG(U)\) commuting with restriction maps: for every
\(V\subset U\),
\[
\begin{tikzcd}
\FF(U) \arrow[r,"\varphi_U"] \arrow[d] & \GG(U) \arrow[d]\\
\FF(V) \arrow[r,"\varphi_V"] & \GG(V).
\end{tikzcd}
\]
The same definition gives morphisms of sheaves.
\end{definition}

\section{Stalks and Germs}

\begin{definition}
Let \(\FF\) be a presheaf on \(X\), and let \(x\in X\).  The \emph{stalk} of
\(\FF\) at \(x\) is
\[
    \FF_x=\varinjlim_{x\in U}\FF(U),
\]
where the limit is taken over open neighbourhoods of \(x\), ordered by reverse
inclusion.  An element of \(\FF_x\) is called a \emph{germ}.  The germ of a
section \(s\in \FF(U)\) at \(x\in U\) is denoted \(s_x\).
\end{definition}

\begin{lemma}
Let \(\FF\) be a presheaf and \(s,t\in \FF(U)\).  Then \(s_x=t_x\) in
\(\FF_x\) for a point \(x\in U\) if and only if there is an open neighbourhood
\(V\subset U\) of \(x\) such that \(s|_V=t|_V\).
\end{lemma}

\begin{proof}
This is the defining equivalence relation in the filtered colimit.  Two pairs
\((U,s)\) and \((U,t)\) represent the same germ at \(x\) precisely when they
have the same image after restriction to some common neighbourhood of \(x\).
\end{proof}

\begin{proposition}
Let \(\FF\) be a sheaf on \(X\), let \(U\subset X\) be open, and let
\(s,t\in \FF(U)\).  If \(s_x=t_x\) for all \(x\in U\), then \(s=t\).
\end{proposition}

\begin{proof}
For each \(x\in U\), the previous lemma gives an open neighbourhood
\(V_x\subset U\) such that \(s|_{V_x}=t|_{V_x}\).  The sets \(V_x\) cover \(U\).
By the uniqueness part of the sheaf axiom, two sections of \(\FF(U)\) with the
same restrictions to a cover are equal.  Hence \(s=t\).
\end{proof}

\begin{proposition}
Let \(\varphi:\FF\to \GG\) be a morphism of sheaves of sets.  If every stalk map
\(\varphi_x:\FF_x\to \GG_x\) is an isomorphism, then \(\varphi\) is an
isomorphism of sheaves.
\end{proposition}

\begin{proof}
It is enough to prove that each map \(\varphi_U:\FF(U)\to\GG(U)\) is bijective.
Injectivity follows from the preceding proposition: if
\(\varphi_U(s)=\varphi_U(t)\), then \(\varphi_x(s_x)=\varphi_x(t_x)\) for all
\(x\in U\), hence \(s_x=t_x\) for all \(x\), and so \(s=t\).

For surjectivity, take \(g\in \GG(U)\).  Since \(\varphi_x\) is surjective, for
each \(x\in U\) there are an open neighbourhood \(U_x\subset U\) and a section
\(f_x\in \FF(U_x)\) such that \(\varphi(f_x)_x=g_x\).  After shrinking \(U_x\),
we may assume \(\varphi(f_x)=g|_{U_x}\).  On overlaps \(U_x\cap U_y\), the
sections \(f_x\) and \(f_y\) have the same image under \(\varphi\).  By
injectivity already proved for the open set \(U_x\cap U_y\), they are equal
there.  The sheaf axiom glues the \(f_x\) to a section \(f\in\FF(U)\), and
\(\varphi(f)=g\).
\end{proof}

\section{Sheafification}

\begin{definition}
Let \(\FF\) be a presheaf of sets on \(X\).  Define a topological space
\[
    E(\FF)=\coprod_{x\in X}\FF_x
\]
with projection \(p:E(\FF)\to X\).  For each section \(s\in \FF(U)\), define
\(\widetilde{s}:U\to E(\FF)\) by \(x\mapsto s_x\).  The topology on \(E(\FF)\)
is generated by the subsets \(\widetilde{s}(U)\).  The \emph{sheafification}
\(\FF^+\) is the sheaf of continuous sections \(U\to E(\FF)\) of \(p\) which
locally are of the form \(\widetilde{s}\).
\end{definition}

\begin{proposition}
For every presheaf \(\FF\), the assignment \(\FF^+\) is a sheaf, and there is a
natural morphism \(a:\FF\to \FF^+\) inducing isomorphisms on all stalks.
\end{proposition}

\begin{proof}
The map \(a_U\) sends \(s\in\FF(U)\) to the local section
\(\widetilde{s}:U\to E(\FF)\).  Its germ at \(x\) is \(s_x\), so \(a_x\) is the
identity map \(\FF_x\to(\FF^+)_x\) once the latter stalk is identified with the
fiber \(p^{-1}(x)=\FF_x\).

We verify the sheaf axiom.  Suppose \(U=\bigcup_i U_i\) and
\(\sigma_i\in\FF^+(U_i)\) agree on overlaps.  Define \(\sigma:U\to E(\FF)\) by
\(\sigma(x)=\sigma_i(x)\) for any \(i\) with \(x\in U_i\).  This is well
defined because the \(\sigma_i\) agree on intersections.  Continuity and the
property of being locally represented by sections of \(\FF\) are local on
\(U\), so \(\sigma\in\FF^+(U)\).  Uniqueness is pointwise.  Hence \(\FF^+\) is a
sheaf.
\end{proof}

\begin{proposition}[Universal property]
Let \(\FF\) be a presheaf and let \(\GG\) be a sheaf on \(X\).  Every morphism
\(\varphi:\FF\to\GG\) factors uniquely through the sheafification:
\[
\begin{tikzcd}
\FF \arrow[r,"a"] \arrow[dr,"\varphi"'] & \FF^+ \arrow[d,dashed,"\varphi^+"]\\
& \GG .
\end{tikzcd}
\]
\end{proposition}

\begin{proof}
For \(\sigma\in\FF^+(U)\), choose an open cover \(U=\bigcup_i U_i\) such that
\(\sigma|_{U_i}=a(s_i)\) for sections \(s_i\in\FF(U_i)\).  Define
\(\varphi^+(\sigma)|_{U_i}=\varphi(s_i)\).  On \(U_i\cap U_j\), the germs of
\(s_i\) and \(s_j\) agree at every point because they define the same section
of \(\FF^+\).  Hence the germs of \(\varphi(s_i)\) and \(\varphi(s_j)\) agree
at every point.  Since \(\GG\) is a sheaf, these two sections agree on the
overlap.  Thus the \(\varphi(s_i)\) glue to a unique section of \(\GG(U)\).
This gives a morphism \(\varphi^+:\FF^+\to\GG\), and by construction
\(\varphi^+\circ a=\varphi\).

If \(\psi:\FF^+\to\GG\) also satisfies \(\psi\circ a=\varphi\), then
\(\psi\) and \(\varphi^+\) agree on the local sections \(a(s)\), and every
section of \(\FF^+\) is locally of this form.  By the sheaf axiom for \(\GG\),
\(\psi=\varphi^+\).
\end{proof}

\section{Kernels, Images, and Cokernels}

\begin{definition}
Let \(\varphi:\FF\to\GG\) be a morphism of sheaves of abelian groups.
The \emph{kernel} sheaf is the presheaf
\[
    U\longmapsto \ker(\FF(U)\to\GG(U)).
\]
This presheaf is already a sheaf.  The \emph{image} sheaf
\(\im(\varphi)\) is the sheafification of the presheaf
\[
    U\longmapsto \im(\FF(U)\to\GG(U)).
\]
The \emph{cokernel} sheaf is the sheafification of
\[
    U\longmapsto \coker(\FF(U)\to\GG(U)).
\]
\end{definition}

\begin{proposition}
Let \(\varphi:\FF\to\GG\) be a morphism of sheaves of abelian groups.  For every
\(x\in X\),
\[
    (\ker\varphi)_x=\ker(\varphi_x),\qquad
    (\im\varphi)_x=\im(\varphi_x),\qquad
    (\coker\varphi)_x=\coker(\varphi_x).
\]
\end{proposition}

\begin{proof}
Filtered colimits of abelian groups are exact.  Since stalks are filtered
colimits over neighbourhoods of \(x\), taking stalks commutes with kernels,
images, and cokernels at the presheaf level.  Sheafification does not change
stalks by the preceding section.  The displayed equalities follow.
\end{proof}

\begin{definition}
A sequence of sheaves of abelian groups
\[
    \FF'\longrightarrow \FF\longrightarrow \FF''
\]
is \emph{exact} if the image sheaf of the first morphism equals the kernel
sheaf of the second.
\end{definition}

\begin{proposition}
A sequence of sheaves of abelian groups is exact if and only if the induced
sequence on stalks is exact at every point.
\end{proposition}

\begin{proof}
By the previous proposition, the stalk of the image sheaf is the image of the
stalk map and the stalk of the kernel sheaf is the kernel of the stalk map.
Thus exactness of sheaves implies exactness on stalks.  Conversely, if the
image and kernel have the same stalks at every point, then the natural inclusion
\(\im\to\ker\) is an isomorphism on stalks.  A morphism of sheaves that is an
isomorphism on stalks is an isomorphism, so \(\im=\ker\).
\end{proof}

\section{Direct and Inverse Image}

\begin{definition}
Let \(f:X\to Y\) be a continuous map.  If \(\FF\) is a sheaf on \(X\), its
\emph{direct image} \(f_*\FF\) is the sheaf on \(Y\) defined by
\[
    (f_*\FF)(V)=\FF(f^{-1}V).
\]
If \(\GG\) is a sheaf on \(Y\), its \emph{inverse image} \(f^{-1}\GG\) is the
sheafification of the presheaf
\[
    U\longmapsto \varinjlim_{V\supset f(U)}\GG(V),
\]
where \(V\) ranges over open subsets of \(Y\) containing \(f(U)\).
\end{definition}

\begin{proposition}
For sheaves of sets, abelian groups, or rings there is a natural adjunction
\[
    \Hom_X(f^{-1}\GG,\FF)\simeq \Hom_Y(\GG,f_*\FF).
\]
\end{proposition}

\begin{proof}
Given \(\alpha:f^{-1}\GG\to\FF\), compose
\(\GG(V)\to(f^{-1}\GG)(f^{-1}V)\) with
\(\alpha_{f^{-1}V}\) to obtain a section of \(\FF(f^{-1}V)\).  This is
functorial in \(V\), hence gives \(\GG\to f_*\FF\).

Conversely, given \(\beta:\GG\to f_*\FF\), an element of the presheaf defining
\(f^{-1}\GG\) over \(U\) is locally represented by a section
\(s\in\GG(V)\) with \(f(U)\subset V\).  The section
\(\beta(s)\in\FF(f^{-1}V)\) restricts to \(\FF(U)\).  These local assignments
are compatible with restriction, so by the universal property of
sheafification they define \(f^{-1}\GG\to\FF\).  The two constructions are
inverse because both are obtained by restricting the same sections.
\end{proof}

\section{Sheaves on Bases}

\begin{definition}
Let \(X\) be a topological space and let \(\mathcal B\) be a basis for its
topology.  A \emph{presheaf on the basis} \(\mathcal B\) assigns an object
\(\FF(U)\) to every \(U\in\mathcal B\), together with restriction maps for
inclusions \(V\subset U\) in \(\mathcal B\), satisfying the usual identity and
transitivity rules.  It is a \emph{sheaf on the basis} if the sheaf gluing
axiom holds whenever \(U\in\mathcal B\) is covered by basis opens
\(\{U_i\}\), with compatibility checked after refining intersections by basis
opens.
\end{definition}

\begin{proposition}
Let \(\mathcal B\) be a basis of \(X\).  Every sheaf \(\FF\) on \(X\) restricts
to a sheaf on \(\mathcal B\).  Conversely, every sheaf \(\FF_{\mathcal B}\) on
\(\mathcal B\) extends uniquely, up to unique isomorphism, to a sheaf \(\FF\) on
\(X\) such that \(\FF(U)=\FF_{\mathcal B}(U)\) for \(U\in\mathcal B\).
\end{proposition}

\begin{proof}
Restriction is immediate.  For the converse, define a presheaf on all open
sets by declaring \(\FF(U)\) to be the set of families
\[
    (s_B)_{B\subset U,\ B\in\mathcal B}
\]
with \(s_B\in\FF_{\mathcal B}(B)\), compatible under restriction:
if \(C\subset B\subset U\), then \(s_B|_C=s_C\).  Restriction from \(U\) to
\(V\) forgets the components not contained in \(V\).  This is a sheaf.  Indeed,
given a cover \(U=\bigcup_i U_i\) and compatible sections over the \(U_i\),
define the component over \(B\subset U\) by choosing a basis cover of \(B\) by
opens contained in some \(U_i\), gluing there by the basis sheaf axiom, and
using uniqueness to see that the result is independent of the chosen
refinement.

If \(U\in\mathcal B\), the map \(\FF_{\mathcal B}(U)\to\FF(U)\) sends a section
to its restrictions to all basis opens \(B\subset U\).  This is injective by
the uniqueness part of the basis sheaf axiom and surjective by gluing the
compatible family over the basis cover of \(U\).  Thus the extension agrees
with \(\FF_{\mathcal B}\) on \(\mathcal B\).  Any other extension has the same
sections on a basis and therefore is uniquely isomorphic to this one.
\end{proof}

\begin{corollary}
On an affine scheme \(X=\Spec A\), an \(\OO_X\)-module is determined by its
values on the principal opens \(D(f)\), provided these values satisfy the
sheaf axiom for principal open covers.
\end{corollary}

\begin{proof}
The principal opens form a basis and finite intersections satisfy
\(D(f)\cap D(g)=D(fg)\).  Apply the previous proposition.
\end{proof}

\section{Local Constructions}

\begin{definition}
Let \(\FF\) be a sheaf of abelian groups on \(X\).  The \emph{support} of
\(\FF\) is
\[
    \Supp(\FF)=\{x\in X:\FF_x\ne0\}.
\]
For a section \(s\in\FF(U)\), the support of \(s\) is
\[
    \Supp(s)=\{x\in U:s_x\ne0\}.
\]
\end{definition}

\begin{proposition}
If \(\FF\) is a sheaf of abelian groups and \(s\in\FF(U)\), then
\(\Supp(s)\) is closed in \(U\) whenever equality with the zero section is an
open condition.  In particular this holds for sheaves of modules on a scheme.
\end{proposition}

\begin{proof}
The complement of \(\Supp(s)\) is the set of \(x\in U\) for which \(s_x=0\).
By the definition of germs, \(s_x=0\) if and only if there is an open
neighbourhood \(V\subset U\) of \(x\) such that \(s|_V=0\).  Thus the complement
is a union of open sets.  For sheaves of modules this condition is always
satisfied because \(0\) is a distinguished section and equality of sections is
local.
\end{proof}

\begin{definition}
Let \(i:Z\hookrightarrow X\) be the inclusion of a subset with the subspace
topology, and let \(A\) be an abelian group.  The \emph{skyscraper sheaf}
\(i_*A\) supported on \(Z\) is the direct image of the constant sheaf \(A_Z\) on
\(Z\).  If \(Z=\{x\}\), we write \(A_x\) for this sheaf.
\end{definition}

\begin{proposition}
Let \(x\in X\) be a closed point and let \(A_x\) be the skyscraper sheaf with
value \(A\).  Then \((A_x)_x=A\) and \((A_x)_y=0\) for \(y\ne x\).
\end{proposition}

\begin{proof}
By definition \(A_x(U)=A\) if \(x\in U\), and \(A_x(U)=0\) if \(x\notin U\),
with the evident restriction maps.  Taking the filtered colimit over
neighbourhoods of \(x\) gives \(A\).  If \(y\ne x\), then \(X\setminus\{x\}\)
is an open neighbourhood of \(y\) on which the sheaf has value \(0\), so the
stalk at \(y\) is \(0\).
\end{proof}

\begin{definition}
Let \(j:U\hookrightarrow X\) be an open immersion and let \(\FF\) be a sheaf on
\(U\).  The \emph{extension by zero} \(j_!\FF\) is the sheaf associated to the
presheaf
\[
    V\longmapsto
    \{s\in \FF(V\cap U):\Supp(s)\text{ is closed in }V\}.
\]
\end{definition}

\begin{proposition}
For an open immersion \(j:U\hookrightarrow X\), the functor \(j_!\) is left
adjoint to restriction:
\[
    \Hom_X(j_!\FF,\GG)\simeq \Hom_U(\FF,\GG|_U).
\]
\end{proposition}

\begin{proof}
A morphism \(j_!\FF\to\GG\), after restricting to \(U\), gives
\(\FF\to\GG|_U\) because \(j_!\FF|_U\simeq\FF\).  Conversely, given
\(\alpha:\FF\to\GG|_U\), define on a section \(s\in j_!\FF(V)\) the section of
\(\GG(V)\) which equals \(\alpha(s)\) on \(V\cap U\) and equals \(0\) outside
the closed support of \(s\).  The closed support condition makes these local
definitions compatible along the open cover \(V=(V\cap U)\cup(V\setminus
\Supp(s))\).  The construction is compatible with restrictions, hence gives a
morphism \(j_!\FF\to\GG\).  The two constructions are inverse.
\end{proof}

\section{Gluing Sheaves}

\begin{proposition}[Gluing sheaves]
Let \(X=\bigcup_i U_i\) be an open cover.  Suppose that for each \(i\) we are
given a sheaf \(\FF_i\) on \(U_i\), and for each pair \(i,j\) an isomorphism
\[
    \varphi_{ij}:\FF_j|_{U_i\cap U_j}\longrightarrow
    \FF_i|_{U_i\cap U_j}
\]
such that \(\varphi_{ii}=\id\) and
\[
    \varphi_{ij}\circ\varphi_{jk}=\varphi_{ik}
    \quad\text{on }U_i\cap U_j\cap U_k.
\]
Then there exists a sheaf \(\FF\) on \(X\), unique up to unique isomorphism,
with isomorphisms \(\FF|_{U_i}\simeq\FF_i\) inducing the given transition
isomorphisms.
\end{proposition}

\begin{proof}
For an open set \(V\subset X\), define \(\FF(V)\) to be the set of families
\[
    (s_i)_i,\qquad s_i\in\FF_i(V\cap U_i),
\]
such that on \(V\cap U_i\cap U_j\) one has
\[
    s_i=\varphi_{ij}(s_j).
\]
Restrictions are defined componentwise.  The sheaf axiom for \(\FF\) follows
from the sheaf axiom for the \(\FF_i\): compatible local families glue
componentwise, and the transition compatibility survives because it can be
checked after restriction to the covering opens.  Restricting \(\FF\) to
\(U_i\) recovers \(\FF_i\), since a compatible family is determined by its
\(i\)-th component and the transition maps.  Uniqueness follows because any
sheaf with these restrictions has the same description by compatible local
sections.
\end{proof}

\begin{proposition}[Gluing morphisms]
Let \(X=\bigcup_i U_i\), and let \(\FF,\GG\) be sheaves on \(X\).  Suppose
morphisms \(\alpha_i:\FF|_{U_i}\to\GG|_{U_i}\) agree on every
\(U_i\cap U_j\).  Then there is a unique morphism
\(\alpha:\FF\to\GG\) whose restriction to \(U_i\) is \(\alpha_i\).
\end{proposition}

\begin{proof}
For \(V\subset X\) and \(s\in\FF(V)\), the sections
\(\alpha_i(s|_{V\cap U_i})\in\GG(V\cap U_i)\) agree on overlaps by hypothesis.
They glue uniquely to a section of \(\GG(V)\).  This assignment is compatible
with restriction and is additive, so it defines a morphism of sheaves.  Any
global morphism restricting to the \(\alpha_i\) must act in this way on every
section, so it is unique.
\end{proof}

\begin{remark}
This gluing proposition is the sheaf-theoretic form of the gluing constructions
for schemes, quasi-coherent sheaves, and line bundles.  In later chapters the
cocycle condition on triple intersections is the mechanism behind transition
functions.
\end{remark}

\section{Exercises Used Later}

\begin{exercise}[Constant sheaves on connected components]\label{ex:constant-sheaf}
Let \(A\) be an abelian group and let \(A_X\) be the sheaf associated to the
constant presheaf on \(X\).  Show that \(A_X(U)\) is the group of locally
constant maps \(U\to A\).  Deduce that if \(U\) has connected components
\(\{U_\lambda\}\), then \(A_X(U)\simeq\prod_\lambda A\).
\end{exercise}

\begin{solution}
The sheafification of the constant presheaf has stalk \(A\) at every point.  A
section over \(U\) is locally represented by a constant element of \(A\), hence
is exactly a locally constant map \(U\to A\).  Conversely any locally constant
map is locally constant with some value \(a\in A\), hence gives a section of
the sheafification.  On a connected component a locally constant map is
constant, and the choices on distinct components are independent, giving the
product decomposition.
\end{solution}

\begin{exercise}[Surjectivity is local on the target]\label{ex:sheaf-surj-local}
Let \(\varphi:\FF\to\GG\) be a morphism of sheaves of abelian groups.  Show that
\(\varphi\) is an epimorphism in the sheaf sense if and only if for every open
\(U\) and every \(t\in\GG(U)\), there is an open cover \(U=\bigcup_i U_i\) and
sections \(s_i\in\FF(U_i)\) with \(\varphi(s_i)=t|_{U_i}\).
\end{exercise}

\begin{solution}
If \(\varphi\) is an epimorphism, then \(\coker\varphi=0\).  Hence for each
\(x\in U\), the germ \(t_x\) lies in the image of \(\FF_x\to\GG_x\).  Thus
there is a neighbourhood \(U_x\) and a section \(s_x\in\FF(U_x)\) mapping to
\(t|_{U_x}\) after possibly shrinking.  These \(U_x\) form the required cover.
Conversely, if such local lifts exist for every \(t\), then every germ of
\(\GG\) lies in the image of the corresponding stalk map.  Therefore the
cokernel has zero stalks, hence is the zero sheaf.
\end{solution}

\begin{exercise}[Exactness after restriction]\label{ex:restriction-exact}
Let \(j:U\hookrightarrow X\) be an open immersion.  Show that restriction
\(\FF\mapsto\FF|_U\) is an exact functor on sheaves of abelian groups.
\end{exercise}

\begin{solution}
The stalk of \(\FF|_U\) at a point \(x\in U\) is \(\FF_x\).  Therefore a
sequence of sheaves on \(X\) remains exact after restriction precisely because
its stalk sequence at each point of \(U\) remains exact.  Exactness of sheaves
is equivalent to exactness on all stalks.
\end{solution}

\begin{exercise}[Skyscraper quotients]\label{ex:skyscraper-quotient}
Let \(X\) be a space, \(x\in X\) a closed point, and \(A\) an abelian group.
Show that any morphism \(\FF\to A_x\) is determined by a homomorphism
\(\FF_x\to A\).
\end{exercise}

\begin{solution}
Taking stalks at \(x\) gives a homomorphism \(\FF_x\to(A_x)_x=A\).  At any
other point the target stalk is zero.  Conversely, given \(u:\FF_x\to A\), a
section \(s\in\FF(U)\) maps to \(u(s_x)\) if \(x\in U\), and to \(0\) if
\(x\notin U\).  This is compatible with restrictions because the point \(x\)
is closed and the only non-zero stalk of \(A_x\) is at \(x\).  The two
constructions are inverse.
\end{solution}
